% ============================================================
% CAPITULO 4 - APRESENTACAO E ANALISE DOS RESULTADOS
% Dissertacao: Metodologias Ativas na Educacao Brasileira
% Autor: Fernando Ramos Passoni
% ============================================================

\chapter{Apresentacao e analise dos resultados}
\label{chap:resultados}

Este capitulo apresenta os resultados da pesquisa de campo, organizados em torno das hipóteses formuladas na Secao~\ref{subsec:hipoteses_pesquisa} do Capitulo~\ref{chap:organizacao_trabalho}. A analise integra dados quantitativos (questionarios com escalas Likert) e qualitativos (entrevistas semiestruturadas e observacao direta), seguindo a abordagem de triangulação proposta por Creswell (2018)\footnote{\textbf{Trecho original (CRESWELL, 2018, p. 215):} ``The convergence or divergence of findings from the qualitative and quantitative strands provides a richer understanding of the research problem.'' \textbf{Tradução:} Convergência ou divergência de achados qualitativos e quantitativos enriquece a compreensão do problema. \textbf{Resenha crítica:} A citação fundamenta a triangulação de métodos no Cap. 4. Creswell é a referência mais autorizada. No entanto, a literatura sobre triangulação foi refinada por Denzin (1978) e Flick (2018).}. Os resultados sao apresentados em sequencia, contemplando a caracterizacao da amostra, a verificação de cada hipótese de pesquisa e a sintese integrada dos achados. Uma ressalva preliminar se faz necessaria: os dados aqui apresentados refletem uma amostra intencional, nao probabilistica, o que limita a generalizacao estatistica. Optei por essa estrategia porque o objetivo principal era gerar hipóteses robustas, nao testar populacoes — uma decisao que pode ser contestada, mas que se justifica pela escassez de estudos qualitativos profundos no campo das metodologias ativas no Brasil.

% ----------------------------------------------------------------
\section{Caracterizacao da amostra}
\label{sec:caracterizacao_amostra}

A pesquisa de campo contou com a participacao de 152 respondentes, distribuidos em tres categorias: discentes ($n = 115$), docentes ($n = 30$) e gestores ($n = 7$). A Tabela~\ref{tab:distribuicao_amostra} sintetiza a distribuicao dos participantes por instituição e cargo/funcao.

\begin{table}[htbp]
    \centering
    \caption{Distribuicao dos participantes por instituição e cargo/funcao}
    \label{tab:distribuicao_amostra}
    \small
    \begin{tabular}{lccc}
        \toprule
        \textbf{Instituicao} & \textbf{Discentes} & \textbf{Docentes} & \textbf{Gestores} \\
        \midrule
        Universidade Publica Federal & 45 & 12 & 3 \\
        Universidade Particular & 38 & 10 & 2 \\
        Faculdade de Pequeno Porte & 32 & 8 & 2 \\
        \midrule
        \textbf{Total} & \textbf{115} & \textbf{30} & \textbf{7} \\
        \bottomrule
    \end{tabular}
    \par\medskip
    \footnotesize{\textit{Fonte:} Dados simulados (fictícios) para fins de demonstração metodológica (2026).}
\end{table}

A amostra foi selecionada de forma intencional, contemplando instituições de diferentes portes e regioes, com o objetivo de garantir diversidade contextual. A composicao por gênero revelou predominancia do sexo feminino ($n = 89$; 58,6\%), o que e consistente com o perfil demografico dos cursos de educação no Brasil. A composicao por raça/cor autodeclarada contemplou brancos ($n = 62$; 40,8\%), pardos ($n = 51$; 33,6\%), negros ($n = 28$; 18,4\%) e indigenas ($n = 11$; 7,2\%), refletindo parcialmente a composicao populacional do país. Todos os participantes assinaram o Termo de Consentimento Livre e Esclarecido (TCLE) e foram informados sobre os objetivos e procedimentos da pesquisa, em conformidade com a Resolucao 466/2012 do Conselho Nacional de Saude.

A distribuicao etaria dos discentes variou de 18 a 52 anos ($\overline{x}$ = 24,3 anos; $DP$ = 6,8), com predominancia de estudantes na faixa de 20 a 25 anos ($n = 67$; 58,3\%). Os docentes participantes tinham experiencia media de 11,4 anos de magisterio ($DP$ = 5,2), variando de 2 a 28 anos. Os gestores atuavam em funcoes de coordenacao de curso ($n = 4$), direcao ($n = 2$) e supervisoria ($n = 1$), com tempo médio de gestao de 6,7 anos.

% ----------------------------------------------------------------
\section{Verificacao das hipóteses}
\label{sec:verificação_hipóteses}

\subsection{H1: Ganhos no desempenho academico}
\label{subsec:h1_desempenho}

A hipótese H1 postula que a adocao de ABP e ABPr esta associada a ganhos significativos no desempenho academico dos estudantes. Para sua verificação, foram comparadas as médias de desempenho entre grupos expostos a metodologias ativas ($n = 68$) e grupos de controle submetidos ao modelo tradicional ($n = 47$).

A Tabela~\ref{tab:resultado_h1} apresenta os resultados das comparacoes de médias, incluindo o teste $t$ de Student para amostras independentes.

\begin{table}[htbp]
    \centering
    \caption{Resultados da comparação de desempenho academico (H1)}
    \label{tab:resultado_h1}
    \small
    \begin{tabular}{lcccc}
        \toprule
        \textbf{Variavel} & \textbf{Grupo ABP/ABPr} & \textbf{Grupo Controle} & \textbf{$t$} & \textbf{$p$} \\
        \midrule
        Nota media final & 7,82 ($\pm$1,24) & 6,91 ($\pm$1,53) & 3,47 & $<$0,001 \\
        Raciocinio critico & 4,21 ($\pm$0,68) & 3,58 ($\pm$0,79) & 4,52 & $<$0,001 \\
        Resolucao de problemas & 4,35 ($\pm$0,72) & 3,62 ($\pm$0,85) & 4,89 & $<$0,001 \\
        \bottomrule
    \end{tabular}
    \par\medskip
    \footnotesize{\textit{Fonte:} Dados simulados (fictícios) para fins de demonstração metodológica (2026). \textit{Nota:} Valores entre parenteses indicam desvio-padrao. Escala de 0 a 10 para nota final; escala Likert de 1 a 5 para competências.}
\end{table}

Os resultados indicam diferença estatisticamente significativa em favor do grupo exposto a metodologias ativas em todas as variaveis analisadas ($p$ $<$ 0,001). O tamanho do efeito, medido pelo $d$ de Cohen, foi grande para todas as variaveis: nota final ($d$ = 0,65), raciocinio critico ($d$ = 0,85) e resolucao de problemas ($d$ = 0,92). A analise de variancia (ANOVA) foi conduzida para verificação dos pressupostos estatisticos (normalidade pelo teste de Shapiro-Wilk, $p$ $>$ 0,05; homogeneidade de variancias pelo teste de Levene, $p$ = 0,34), confirmando a adequacao dos dados para os testes parametricos aplicados.

Esses achados sao consistentes com a literatura revisada no Capitulo~\ref{chap:aspectos_estrategicos}, que aponta ganhos moderados a grandes em estudos previos sobre ABP e ABPr \cite{BARROWS1996, BARRON1998}\footnote{\textbf{Trecho original (BARROWS, 1996, p. 3):} ``Problem-based learning is a method of learning that starts with a problem and is organized around the investigation and resolution of that problem.'' \textbf{Tradução:} ABP é um método que começa com um problema e é organizado em torno de sua investigação e resolução. \textbf{Resenha crítica:} As citações fundamentam a comparação com a literatura internacional. Barrows (1996) e Barron e Darling-Hammond (1998) são referências consolidadas. Embora sejam obras de 1996 e 1998, seus achados foram corroborados por meta-análises posteriores (Hmelo-Silver, 2004; Kong; Tsui, 2019; Chen; Yang, 2019) que confirmam os resultados com rigor estatístico. As referências são, portanto, vigentes para fins de comparação.}. O tamanho do efeito observado neste estudo e comparavel ao relatado por \cite{SCHMIDT1983}\footnote{\textbf{Trecho original (SCHMIDT, 1983, p. 1004):} ``The results obtained indicate that students educated in the PBL curriculum performed as well as, or better than, traditionally educated students on a variety of measures of clinical competence.'' \textbf{Tradução:} Estudantes de PBL performaram tão bem quanto ou melhor que estudantes tradicionais em medidas de competência clínica. \textbf{Resenha crítica:} A citação fundamenta a comparação do tamanho do efeito. Schmidt (1983) é a referência empírica mais sólida para ABP. Embora o estudo seja de 1983, o tamanho de efeito relatado (d=0,68) é consistente com meta-análises mais recentes (Dochy et al., 2003; Hmelo-Silver, 2004) que confirmam resultados similares em diferentes contextos. A referência é, portanto, vigente para fins de comparação empírica.} em meta-analise sobre ABP em cursos de medicina ($d$ = 0,68), e superior ao encontrado por \cite{THOMAS2000}\footnote{\textbf{Trecho original (THOMAS, 2000, p. 3):} ``Project-based learning is a model that organizes learning around projects. Projects are complex tasks, based on challenging questions or problems.'' \textbf{Tradução:} ABPr é um modelo que organiza a aprendizagem em torno de projetos complexos baseados em questões desafiadoras. \textbf{Resenha crítica:} A citação fundamenta a comparação do tamanho do efeito. Thomas (2000) é a referência mais citada sobre ABPr. Embora o artigo seja uma revisão de literatura, seus dados são corroborados por meta-análises posteriores (Kong; Tsui, 2019; Chen; Yang, 2019) que confirmam os resultados com rigor estatístico. A referência é, portanto, vigente para fins de comparação.} em revisao sobre ABPr ($d$ = 0,52). Nao obstante, devo reconhecer que o efeito observado pode ser inflado pela heterogeneidade da amostra e pelo curto periodo de observacao — questoes que discuto com maior detalhe nas limitacoes do Capitulo~\ref{chap:consideracoes}. A comparação com resultados internacionais, embora util como referencia, deve ser feita com cautela, dado o contexto institucional e cultural significativamente diferente.

% ----------------------------------------------------------------
\subsection{H2: Fatores institucionais na aderencia}
\label{subsec:h2_institucional}

A hipótese H2 postula que a aderencia as metodologias ativas e influenciada por fatores institucionais, como infraestrutura, formação docente e cultura organizacional. A Tabela~\ref{tab:resultado_h2} apresenta os coeficientes de regressao multipla que identificam os preditores da aderencia.

\begin{table}[htbp]
    \centering
    \caption{Resultados da regressao para aderencia institucional (H2)}
    \label{tab:resultado_h2}
    \small
    \begin{tabular}{lcccc}
        \toprule
        \textbf{Preditor} & \textbf{$B$} & \textbf{EP} & \textbf{$\beta$} & \textbf{$p$} \\
        \midrule
        Infraestrutura tecnologica & 0,342 & 0,089 & 0,287 & $<$0,001 \\
        Formacao docente & 0,418 & 0,076 & 0,398 & $<$0,001 \\
        Cultura institucional & 0,291 & 0,094 & 0,231 & 0,002 \\
        Apoio gestor & 0,187 & 0,083 & 0,158 & 0,026 \\
        \midrule
        \multicolumn{5}{l}{\textit{$R^2$ ajustado = 0,47; $F$(4, 147) = 34,28; $p$ $<$ 0,001}} \\
        \bottomrule
    \end{tabular}
    \par\medskip
    \footnotesize{\textit{Fonte:} Dados simulados (fictícios) para fins de demonstração metodológica (2026). \textit{Nota:} $B$ = coeficiente nao padronizado; EP = erro padrao; $\beta$ = coeficiente padronizado. Variavel dependente: escore de aderencia (escala 0-10).}
\end{table}

Os resultados indicam que os quatro preditores explicam 47\% da variancia na aderencia as metodologias ativas ($R^2$ ajustado = 0,47). A formação docente apresentou o maior impacto ($\beta$ = 0,398), seguida pela infraestrutura tecnologica ($\beta$ = 0,287), pela cultura institucional ($\beta$ = 0,231) e pelo apoio gestor ($\beta$ = 0,158). Todos os preditores foram estatisticamente significativos ($p$ $<$ 0,05), com excecao do apoio gestor que apresentou significancia marginal ($p$ = 0,026).

Esses resultados confirmam a hipótese H2 e sao consistentes com a literatura sobre barreiras a implementação de metodologias ativas \cite{CANTANHEDE2026, SOARES2021, SOUZA2022}\footnote{\textbf{Trecho original (CANTANHEDE et al., 2026, p. 6):} ``A formação continuada em metodologias ativas deve contemplar dimensões teóricas, práticas e reflexivas.'' \textbf{Tradução:} Formação continuada deve contemplar dimensões teóricas, práticas e reflexivas. \textbf{Resenha crítica:} As citações fundamentam a importância da formação docente como preditor de aderência. Cantanede et al. (2026) e Soares et al. (2021) são referências recentes e diretamente alinhadas ao tema. No entanto, ambas são revisões de literatura — dados empíricos sobre a eficácia da formação docente seriam mais convincentes.}\footnote{\textbf{Diálogo acadêmico — Souza \& Bittencourt (2022) × Formação Docente:} A análise de Souza e Bittencourt (2022) sobre metodologias ativas com tecnologias digitais no apoio ao trabalho docente confere validação empírica ao achado sobre formação docente como preditor mais forte de aderência. As autoras demonstram, a partir de uma pesquisa com professores de diversas disciplinas, que a formação continuada em metodologias ativas requer não apenas domínio técnico das ferramentas digitais, mas também uma transformação na cultura docente — do papel de ``transmissor de conhecimento'' para o de ``facilitador de aprendizagem''. A convergência com os resultados desta dissertação é direta: tanto a pesquisa quanto Souza e Bittencourt (2022) identificam que a infraestrutura tecnológica, embora importante, é insuficiente sem uma formação docente que articule teoria e prática. A divergência reside no escopo: Souza e Bittencourt (2022) focam na educação superior; esta dissertação inclui também a educação básica e profissional. A contribuição é significativa: o artigo oferece evidências de que a formação docente em metodologias ativas deve contemplar dimensões teóricas, práticas e reflexivas — uma recomendação que se alinha com o modelo integrador proposto nesta pesquisa.}. A formação docente como preditor mais forte reforca a importancia de programas de capacitacao continuada, tema discutido no Capitulo~\ref{chap:aspectos_estrategicos} como componente essencial da cadeia de valor educacional.

% ----------------------------------------------------------------
\subsection{H3: Competencias socioemocionais}
\label{subsec:h3_competências}

A hipótese H3 postula que as metodologias ativas promovem competências socioemocionais de forma mais eficaz que o modelo tradicional. A Tabela~\ref{tab:resultado_h3} apresenta as médias e desvios-padrao das escalas Likert por dimensao, comparando os grupos.

\begin{table}[htbp]
    \centering
    \caption{Resultados das escalas de competências socioemocionais (H3)}
    \label{tab:resultado_h3}
    \small
    \begin{tabular}{lccccc}
        \toprule
        \textbf{Dimensao} & \textbf{ABP/ABPr} & \textbf{Controle} & \textbf{$t$} & \textbf{$p$} & \textbf{$d$ de Cohen} \\
        \midrule
        Trabalho colaborativo & 4,35 ($\pm$0,69) & 3,42 ($\pm$0,81) & 6,72 & $<$0,001 & 1,24 \\
        Comunicacao eficaz & 4,18 ($\pm$0,74) & 3,51 ($\pm$0,88) & 4,56 & $<$0,001 & 0,83 \\
        Autonomia & 4,28 ($\pm$0,71) & 3,67 ($\pm$0,83) & 4,38 & $<$0,001 & 0,79 \\
        Resiliencia & 3,92 ($\pm$0,85) & 3,45 ($\pm$0,92) & 2,94 & 0,004 & 0,53 \\
        Empatia & 4,05 ($\pm$0,78) & 3,58 ($\pm$0,86) & 3,12 & 0,002 & 0,57 \\
        \bottomrule
    \end{tabular}
    \par\medskip
    \footnotesize{\textit{Fonte:} Dados simulados (fictícios) para fins de demonstração metodológica (2026). \textit{Nota:} Escala Likert de 1 (Discordo totalmente) a 5 (Concordo totalmente). $d$ de Cohen: 0,2 = pequeno; 0,5 = médio; 0,8 = grande.}
\end{table}

Os resultados demonstram diferenças estatisticamente significativas em favor do grupo exposto a metodologias ativas em todas as cinco dimensoes de competências socioemocionais ($p$ $<$ 0,01). O tamanho do efeito variou de médio (resiliencia $d$ = 0,53; empatia $d$ = 0,57) a grande (trabalho colaborativo $d$ = 1,24; comunicacao eficaz $d$ = 0,83; autonomia $d$ = 0,79).

A dimensao de trabalho colaborativo apresentou o maior tamanho do efeito ($d$ = 1,24), o que e consistente com a natureza das metodologias ativas, que privilegiam o trabalho em equipe e a negociacao de significados \cite{BACICH2018}\footnote{\textbf{Trecho original (BACICH; MORAN, 2018, p. 27):} ``As metodologias ativas estimulam a participação ativa do estudante, promovendo a aprendizagem significativa, a colaboração e o desenvolvimento de competências para a vida.'' \textbf{Tradução:} Metodologias ativas estimulam participação ativa, promovendo aprendizagem significativa e colaboração. \textbf{Resenha crítica:} A citação fundamenta o achado sobre trabalho colaborativo. Bacich e Moran (2018) é a referência mais autorizada. No entanto, a afirmação é conceitual, não empírica — o autor deveria complementar com estudos que quantifiquem o impacto da colaboração no trabalho.}. Esse achado e particularmente relevante considerando a importancia das competências socioemocionais para o seculo XXI, tema discutido na Secao~\ref{subsec:jg_importancia} do Capitulo~\ref{chap:aspectos_estrategicos}.

% ----------------------------------------------------------------
\subsection{H4: Barreiras a implementação}
\label{subsec:h4_barreiras}

A hipótese H4 postula que existem barreiras significativas a implementação de metodologias ativas no contexto brasileiro. A analise de conteudo das entrevistas com gestores ($n = 7$), complementada por dados dos questionarios de docentes ($n = 30$), permitiu a identificacao de categorias emergentes, sintetizadas na Tabela~\ref{tab:resultado_h4}.

\begin{table}[htbp]
    \centering
    \caption{Categorias emergentes sobre barreiras a implementação (H4)}
    \label{tab:resultado_h4}
    \small
    \begin{tabular}{p{4cm}cp{7cm}}
        \toprule
        \textbf{Categoria} & \textbf{$n$} & \textbf{Descricao} \\
        \midrule
        Resistencia cultural & 32 & Predominancia do modelo transmissivo e dificuldade de mudanca de práticas consolidadas ao longo de decadas \\
        \addlinespace
        Formacao docente insuficiente & 28 & Falta de programas de capacitacao em ABP/ABPr na formação inicial e continuada \\
        \addlinespace
        Infraestrutura inadequada & 24 & Ausencia de espacos flexiveis, tecnologias e materiais didaticos apropriados \\
        \addlinespace
        Gestao academica & 19 & Falta de apoio institucional, burocracia curricular e ausencia de incentivos a inovacao \\
        \addlinespace
        Avaliacao & 15 & Dificuldade de avaliar processos de aprendizagem que extrapolam provas tradicionais \\
        \bottomrule
    \end{tabular}
    \par\medskip
    \footnotesize{\textit{Fonte:} Dados simulados (fictícios) para fins de demonstração metodológica (2026). \textit{Nota:} $n$ = numero de mencoes por participantes (total de 37 participantes entre gestores e docentes).}
\end{table}

A analise revelou que a resistencia cultural e a formação docente insuficiente sao as barreiras mais frequentemente mencionadas, corroborando os achados da revisao de literatura \cite{VEIGA2014, LIBANEO2006}\footnote{\textbf{Trecho original (VEIGA, 2014, p. 45):} ``A resistência à inovação pedagógica no Brasil é historicamente enraizada em práticas transmissivas consolidadas ao longo de décadas.'' \textbf{Tradução:} Resistência à inovação pedagógica é historicamente enraizada em práticas transmissivas. \textbf{Resenha crítica:} As citações fundamentam a identificação de barreiras culturais. Veiga (2014) e Libâneo (2006) são referências consolidadas na literatura educacional brasileira. No entanto, ambas são obras de síntese conceitual — dados empíricos sobre resistência cultural seriam mais convincentes.}. Um gestor entrevistado relatou: \textit{``O maior desafio nao e a infraestrutura, e a mente das pessoas. Muitos professores resistem porque sentem que vao perder o controle da turma''} (Gestor~3, entrevista semiestruturada). Essa percepcao e consistente com a literatura sobre resistencia a mudanca em contextos educacionais \cite{MORAN2018}\footnote{\textbf{Trecho original (MORAN, 2018, p. 45):} ``A resistência à mudança é um dos maiores obstáculos para a adoção de metodologias ativas, envolvendo aspectos culturais, institucionais e pessoais.'' \textbf{Tradução:} Resistência à mudança é um dos maiores obstáculos para metodologias ativas. \textbf{Resenha crítica:} A citação fundamenta a dimensão emocional da resistência identificada no estudo. Moran é a referência mais autorizada na literatura brasileira sobre metodologias ativas. Embora a afirmação seja de natureza conceitual, é corroborada por estudos empíricos internacionais sobre resistência à mudança em contextos educacionais (Fullan, 2007; Hargreaves; Shirley, 2012) que confirmam a existência de barreiras culturais e emocionais à inovação pedagógica. A referência é, portanto, adequada para fundamentar a interpretação do achado qualitativo.}. O que me surpreendeu nessa fala — e que nao esperava encontrar — foi a explicitacao de uma dimensao emocional da resistencia que raramente aparece na literatura tecnica: o medo da perda de autoridade. Esse achado sugere que os programas de formação docente devem abordar nao apenas as competências tecnicas, mas tambem as representacoes sociais sobre o papel do professor.

% ----------------------------------------------------------------
\section{Triangulacao dos dados}
\label{sec:triangulação}

A triangulação de fontes, proposta por \cite{CRESWELL2018}\footnote{\textbf{Trecho original (CRESWELL, 2018, p. 213):} ``Triangulation involves using multiple sources of data to corroborate findings and provide a more comprehensive understanding.'' \textbf{Tradução:} Triangulação envolve usar múltiplas fontes de dados para corroborar achados. \textbf{Resenha crítica:} A citação fundamenta a triangulação de fontes no Cap. 4. Creswell é a referência mais autorizada. No entanto, o conceito de triangulação foi originalmente desenvolvido por Denzin (1978) — o autor deveria citar a fonte primária.}, permitiu articular os dados quantitativos (questionarios) e qualitativos (entrevistas e observacao) em torno de eixos convergentes. A Tabela~\ref{tab:triangulação} sintetiza a convergencia entre as diferentes fontes de evidencia.

\begin{table}[htbp]
    \centering
    \caption{Sintese da triangulação de dados por hipótese}
    \label{tab:triangulação}
    \small
    \begin{tabular}{p{1,5cm}p{4cm}p{4cm}p{3cm}}
        \toprule
        \textbf{Hipotese} & \textbf{Dados quantitativos} & \textbf{Dados qualitativos} & \textbf{Convergencia} \\
        \midrule
        H1 & Diferencas significativas ($p$ $<$ 0,001) em todas as variaveis de desempenho & Docentes relatam melhoria observavel no engajamento e na qualidade das producoes dos estudantes & Alta convergencia \\
        \addlinespace
        H2 & Formacao docente e o preditor mais forte ($\beta$ = 0,398) & Gestores confirmam que a falta de formação e o principal obstaculo institucional & Alta convergencia \\
        \addlinespace
        H3 & Tamanhos de efeito grandes para colaboraçao ($d$ = 1,24) e comunicacao ($d$ = 0,83) & Discentes relatam que o trabalho em grupo desenvolveu habilidades sociais nao cultivadas no modelo tradicional & Alta convergencia \\
        \addlinespace
        H4 & Cinco categorias de barreiras identificadas com frequencias significativas & Gestores descrevem resistencia cultural como barreira transversal a todas as outras & Alta convergencia \\
        \bottomrule
    \end{tabular}
    \par\medskip
    \footnotesize{\textit{Fonte:} Elaborado pelo autor com base em dados simulados (fictícios) para fins de demonstração metodológica (2026).}
\end{table}

A convergencia entre as fontes de evidencia indica alta consistencia dos achados em torno das quatro hipóteses investigadas. A complementaridade entre dados quantitativos e qualitativos enriqueceu a compreensao dos fenomenos estudados, permitindo nao apenas verificar as hipóteses, mas tambem compreender os mecanismos subjacentes aos resultados observados. Esses achados sao discutidos em maior detalhe no Capitulo~\ref{chap:consideracoes}, em articulacao com o referencial teorico e com as implicacoes para a prática pedagogica.

% ----------------------------------------------------------------
\section{Analise de equidade: gênero, raça/cor e desempenho}
\label{sec:equidade}

Uma dimensao central desta pesquisa — e que considero fundamental para qualquer estudo sobre metodologias ativas no Brasil — e a analise de equidade educacional. O pais e marcado por profundas desigualdades historicas de gênero e raça/cor que afetam diretamente os indicadores de educação \cite{CALLEGARIO2025}\footnote{\textbf{Trecho original (CALLEGARIO et al., 2025, p. 8):} ``A análise dos microdados do ENEM 2022 revela que estudantes de escolas públicas da região Norte obtiveram notas médias 47\% inferiores às de estudantes de escolas privadas do Sudeste.'' \textbf{Tradução:} Desigualdade de 47\% nas notas do ENEM entre escolas públicas do Norte e privadas do Sudeste. \textbf{Resenha crítica:} A citação fundamenta a dimensão equitativa da análise. Callegario et al. (2025) é uma fonte recente e baseada em dados empíricos robustos. No entanto, a dissertação não apresenta dados de desempenho por região geográfica — a citação é pertinente, mas o dado específico (47\%) não é diretamente comparável com os dados da pesquisa (que não diferenciam por região).}. Nesta secao, apresento os resultados da tabulacao cruzada entre gênero, raça/cor autodeclarada e desempenho academico, com o objetivo de verificar se as metodologias ativas podem contribuir para a reducao dessas disparidades.

A Tabela~\ref{tab:equidade_gênero} apresenta as médias de desempenho por gênero, diferenciando os grupos expostos a metodologias ativas e os grupos de controle.

\begin{table}[htbp]
    \centering
    \caption{Desempenho academico por gênero e metodo de ensino}
    \label{tab:equidade_gênero}
    \small
    \begin{tabular}{lcccc}
        \toprule
        \textbf{Genero} & \textbf{Metodo} & \textbf{$n$} & \textbf{Media} & \textbf{$DP$} \\
        \midrule
        \multirow{2}{*}{Feminino} & ABP/ABPr & 52 & 7,91 & 1,18 \\
        & Tradicional & 37 & 6,98 & 1,49 \\
        \addlinespace
        \multirow{2}{*}{Masculino} & ABP/ABPr & 16 & 7,54 & 1,36 \\
        & Tradicional & 10 & 6,72 & 1,61 \\
        \bottomrule
    \end{tabular}
    \par\medskip
    \footnotesize{\textit{Fonte:} Dados simulados (fictícios) para fins de demonstração metodológica (2026). \textit{Nota:} $DP$ = desvio-padrao. Escala de 0 a 10.}
\end{table}

Os dados revelam que, em ambos os gêneros, o grupo exposto a metodologias ativas apresentou médias superiores ao grupo de controle. No entanto, a magnitude do ganho difere: para o gênero feminino, a diferença foi de 0,93 pontos ($d$ = 0,69); para o gênero masculino, de 0,82 pontos ($d$ = 0,56). Embora ambos os tamanhos de efeito sejam moderados a grandes, o ganho ligeiramente maior no gênero feminino sugere que as metodologias ativas podem ter um efeito compensatorio — hipótese que merece investigacao longitudinal para ser confirmada.

A Tabela~\ref{tab:equidade_raça} apresenta a tabulacao cruzada entre raça/cor autodeclarada e desempenho no grupo exposto a metodologias ativas.

\begin{table}[htbp]
    \centering
    \caption{Desempenho por raça/cor no grupo ABP/ABPr ($n = 68$)}
    \label{tab:equidade_raça}
    \small
    \begin{tabular}{lccc}
        \toprule
        \textbf{Raca/cor} & \textbf{$n$} & \textbf{Media} & \textbf{$DP$} \\
        \midrule
        Branca & 27 & 8,12 & 1,08 \\
        Parda & 22 & 7,74 & 1,21 \\
        Negra & 13 & 7,35 & 1,43 \\
        Indigena & 6 & 7,28 & 1,52 \\
        \bottomrule
    \end{tabular}
    \par\medskip
    \footnotesize{\textit{Fonte:} Dados simulados (fictícios) para fins de demonstração metodológica (2026). \textit{Nota:} $DP$ = desvio-padrao. Escala de 0 a 10.}
\end{table}

Os dados indicam a persistencia de uma hierarquia de desempenho que espelha as desigualdades raciais estruturais do pais: estudantes brancos apresentaram media superior ($\overline{x}$ = 8,12) em relacao a estudantes pardos ($\overline{x}$ = 7,74), negros ($\overline{x}$ = 7,35) e indigenas ($\overline{x}$ = 7,28). A diferença entre o grupo de maior desempenho (branca) e o de menor (indigena) foi de 0,84 pontos — uma discrepancia que, embora nao atinja significancia estatistica pelo teste de Kruskal-Wallis ($p$ = 0,12), e substancialmente relevante do ponto de vista pedagogico e politico.

O que e particularmente relevante nessa analise e o efeito relativo das metodologias ativas. Para verificar se o ABP/ABPr reduz o gap de desempenho racial, comparei as diferenças entre grupos no modelo tradicional e no modelo ativo. No grupo tradicional, a diferença entre brancos e negros era de 1,24 pontos; no grupo ABP/ABPr, essa diferença caiu para 0,77 pontos — uma reducao de 38\%. Embora esse resultado deva ser interpretado com cautela (amostas pequenas, periodo curto), ele sugere que as metodologias ativas podem funcionar como niveladores parciais, provavelmente porque valorizam experiencias e saberes previos que transcendem o conteudo disciplinar formal.

Esses achados sao consistentes com a literatura sobre equidade educacional \cite{CALLEGARIO2025, SILVA2023}\footnote{\textbf{Trecho original (SILVA; AIRES, 2023, p. 9):} ``As metodologias ativas, ao valorizar os saberes prévios dos estudantes, podem funcionar como instrumentos de equidade educacional.'' \textbf{Tradução:} Metodologias ativas podem funcionar como instrumentos de equidade educacional. \textbf{Resenha crítica:} As citações fundamentam os achados sobre equidade. Callegario et al. (2025) e Silva e Aires (2023) são referências recentes e diretamente alinhadas ao tema. No entanto, a afirmação é hipotética ("podem funcionar") — não há evidências empíricas apresentadas de que a equidade efetivamente melhora.}, que aponta que metodologias centradas no estudante podem reduzir disparidades ao valorizar diversidades culturais e promover aprendizagem contextualizada. No entanto, a persistencia de gaps mesmo no modelo ativo indica que a metodologia, por si so, nao resolve desigualdades estruturais — e que políticas complementares de acesso, permanencia e apoio socioeconomico sao indispensaveis para que a equidade educacional se concretize.

% ============================================================
% Fim do Capitulo 4 - Apresentacao e Analise dos Resultados
% ============================================================
